\documentclass[noauthor,noinstructornotes]{ximera}

\title{HVAC Math}
\author{Math 1193 Design Team}

\begin{document}
\begin{abstract}
    Writing inequalities that describe the behavior of an HVAC system.
\end{abstract}
\maketitle

The heating comes on when the temperature drops strictly below 60 
degrees and the air conditioning comes on when the temperature is 80 
degrees or greater. 

\begin{enumerate}
    \item Write the set of temperatures in interval notation for which
    the heating comes on.
    \item Write the set of temperatures in interval notation for which
    the air conditioning comes on.
    \item Write the set of temperatures in interval notation for which
    the heating \textbf{or} the air conditioning comes on.
    \item Write the set of temperatures in interval notation for which
    the heating \textbf{and} the air conditioning comes on.
\end{enumerate}

\begin{instructorNotes}
This is meant to illustrate the difference between unions and intersections,
as well as interval notation. The terms ``union'' and ``intersection''
aren't necessary to use, however.
\end{instructorNotes}

\end{document}
